% Viabilidad del proyecto. Se retiró de la justificación de
% problema-objetivos.tex porque irá en el documento de metodología (aún no
% existe). Por ahora solo la carga justificacion.tex.
% Redacción en tercera persona.
\ifdefined\justsec\justsec{Viabilidad}\else\section{Viabilidad}\fi
El proyecto propone construir un prototipo de torniquete con un lector
ESP32-CAM que valide un código QR dinámico mostrado en el celular del
estudiante y registre cada ingreso en una base de datos consultable desde
un panel web. El prototipo se probará en un acceso peatonal piloto del
turno vespertino para evaluar su efecto en el tiempo de validación y en la
longitud de la fila.

La viabilidad técnica se sustenta en que los componentes del prototipo, el
ESP32-CAM, el controlador del torniquete, el servidor y la base de datos,
están disponibles comercialmente y corresponden a conocimientos que se han
adquirido en Ingeniería en Sistemas Computacionales. Ya se ha trabajado con
captura de imágenes en pruebas de banco durante la exploración de
alternativas, lo que reduce el riesgo de la etapa de desarrollo. Frente a
NFC, una comparación en transacciones móviles encontró que el QR se procesa
con mayor rapidez, aunque no incorpora cifrado propio
\parencite{harnaningrum2024comparison}; esa carencia se compensa con la firma
HMAC del código.

En cuanto a los recursos económicos, el bajo costo de los componentes
permite construir y probar el prototipo con recursos propios, sin depender
de una inversión institucional. En cuanto al tiempo, el trabajo se organiza
en cinco fases (diagnóstico, diseño de hardware y firmware, desarrollo del
servidor, piloto controlado y evaluación), y el estudio se limita a un solo
acceso peatonal y a la franja de mayor afluencia del turno vespertino, lo
que lo hace abarcable en el periodo escolar.

El acceso a la información y a los participantes es posible porque el
proyecto se desarrolla dentro del \tescha{}: se puede observar el acceso
piloto y aplicar la encuesta a estudiantes y las entrevistas al personal de
vigilancia y de Control Escolar. La instalación del prototipo en el acceso
requiere la autorización de las autoridades del plantel, que se solicitará
antes del piloto.

Queda por confirmar la viabilidad desde el lado del estudiante, es decir, si
cuenta con un celular propio, con conexión a internet y con batería
suficiente al llegar al plantel. Esta condición se indagará con la encuesta
del diagnóstico, y sus resultados indicarán si el sistema puede usarse de
forma general o si el procedimiento manual debe mantenerse para una parte de
los estudiantes.
